\fancyhead[L]{\small KLÜBER ON THE MANNHEIM OBSERVATORY.}
\hypertarget{native-541-lower}{}\label{native-541-lower}
\begin{center}\small Göttingische gelehrte Anzeigen. Nos. 192 and 193. pp. 1913..1916. December 2, 1811.\\[3pt]\rule{.72\linewidth}{.3pt}\end{center}

\emph{The Mannheim Observatory, Described by Its Curator, the Councillor of State and Cabinet Councillor} \textsc{Klüber}. With a lithographic illustration of the observatory. 1811. 62 pages in small folio. Distributed on commission by \textsc{Gottlieb Braun}. Mannheim and Heidelberg.

{\rightskip=0pt plus 1em
Of all the observatories in Germany, Mannheim is now undoubtedly the best equipped with instruments. Everyone who cares about astronomy takes a lively, warm interest in it, so a detailed account of its origins, development and present state can only be welcome. It is doubly pleasing that the report comes from the pen of the observatory's curator, who has its reputation and welfare so much at heart and has already shown this so clearly in his actions.\par}

The foundations of the building, which is solidly built throughout, were laid in 1772 under \textsc{Carl Theodor}, Elector Palatine, according to the plan of the well-known astronomer P. \textsc{Mayer}. No expense was spared to make the building thoroughly durable and handsome. The cost came to over 70000 guilders. This will come as no surprise when one considers the building's considerable height of 111 feet. It was thought necessary at the time, though given present methods of observation it must be regarded as a serious flaw in the design, since it makes observations harder to carry out and reduces their reliability to the same extent. The observatory was equipped with instruments from the finest English instrument makers. An excellent eight-foot mural quadrant by \textsc{Bird} was installed as early as 1775. In 1778 the observatory received a twelve-foot zenith sector by \textsc{Sisson} and a regulator by \textsc{Arnold}. Later \textsc{Ramsden} supplied the six-foot transit telescope (for $145\tfrac{1}{2}$ guineas). A three-foot repeating circle by \textsc{Reichenbach} in Munich was expected and, according to private reports, has in fact already arrived. The first astronomer, \textsc{Christian Mayer}, died in 1783 (his assistant, \textsc{Joh. Metzger}, who was known\SourceLeaf{542} for his special tables for aberration and nutation, died in 1780). From 1784 to 1786 \textsc{Carl König} headed the observatory, and in 1786 and 1787 \textsc{Joh. Nepomuk Fischer} did so. Fischer's designated successor, \textsc{Peter Ungeschick}, died before taking up the post. Finally, since 1788 \textsc{Roger Barry} has been the observatory's astronomer (for a time jointly with \textsc{Henry}). The war greatly hampered the work, and at the end of 1794 all the instruments had to be taken down and packed in crates, where they stayed for 6 years Several times the observatory was even hit by howitzer fire during the bombardment of Mannheim, though without suffering serious damage. The astronomer was subjected to personal mistreatment and persecution, which he could escape only by fleeing. After the Peace of Lunéville, when Mannheim passed to the House of Baden, several scientific and art collections were taken from Mannheim to Munich. The valuable astronomical instruments, however, stayed in Mannheim and were set up again in 1801. Funds for the most necessary upkeep of the observatory and the astronomer's salary were allocated only later.

In 1808 the observatory was placed under the more direct care of Councillor of State \textsc{Klüber}, who did everything he could to revive it. Funds were approved for the upkeep, beautification and furnishing of the building so that it could also occasionally accommodate visiting astronomers, and for lighting, heating, correspondence, a library and other minor needs. A handsome obelisk for the Mire méridienne was put up on the north side, and a second is now being built on the south side. The government has allocated funds to print the large star catalogue on which \textsc{Barry} has been working for many years. A repeating circle by \textsc{Reichenbach} was acquired only recently, as mentioned above. The astronomer himself received a substantial rise in salary and was also assigned a paid attendant. An assistant for the astronomer is now being considered.

The government of Baden and the worthy curator of the observatory certainly deserve the thanks of all friends of this science for this generous support of a fine institution devoted to practical astronomy. One may rightly hope for rich fruits from its active and discerning use.

The rest of the work gives an overview of the determinations of the observatory's position to date, fixing its longitude from Paris at $24^{\mathrm m}31^{\mathrm s}8$ and its latitude at $49^\circ\,29'\,16''$. It also lists all the instruments on hand, gives biographical information about the astronomers who have served at the observatory and provides a complete list of their writings.
